\documentclass[10pt,a4paper]{amsart}
\usepackage[margin=19mm]{geometry}
\usepackage{amsmath,amssymb,mathtools}
\usepackage[scr=rsfs,cal=euler]{mathalfa}
\usepackage{xfrac}
\usepackage[unicode,hidelinks,bookmarks=false]{hyperref}
\newcommand{\ie}{i.e.\ }
\newcommand{\cf}{cf.\ }
\newcommand{\resp}{respectively\ }
\newcommand{\Subsection}{\subsection}
\providecommand{\nobreakdash}{\nobreak\hbox{-}}
\setlength{\emergencystretch}{2em}
\multlinegap0pt
\usepackage{mathtools}
\usepackage[shortlabels]{enumitem}
\usepackage{amssymb}
\usepackage{algorithm}\usepackage{algpseudocode}
\usepackage{float}
\usepackage{bm}
\usepackage{xcolor}
\newtheorem{lemma}{Lemma}
\newcommand{\HH}{\vect{\mathsf{H}}}
\newcommand{\Id}{\operatorname{Id}}
\newcommand{\Jact}{j_{t}}
\newcommand{\aaa}{{\mathsf{a}}}
\newcommand{\aat}{{\mathsf{a}}_t}
\newcommand{\cv}{\overline{v}}
\newcommand{\intG}[1]{\int_{\varGamma}{#1}{\, {d}{s}}}
\newcommand{\intO}[1]{\int_{\varOmega}{#1}{\, {d}{x}}} 
\newcommand{\intS}[1]{\int_{\varSigma}{#1}{\, {d}{s}}} 
\newcommand{\intervalI}{I}
\newcommand{\vect}[1]{\boldsymbol{\mathbf{#1}}}
\newcommand{\vertiii}[1]{{\left\vert\kern-0.25ex\left\vert\kern-0.25ex\left\vert #1 
    \right\vert\kern-0.25ex\right\vert\kern-0.25ex\right\vert}}
\title{Cavity shape reconstruction with a homogeneous Robin condition via a constrained coupled complex boundary method with ADMM}
\author{Mustapha Essahraoui, El Mehdi Cherrat, Lekbir Afraites, Julius Fergy Tiongson Rabago}
\date{Source: arXiv:2605.12202v2 (CC BY 4.0)}
\begin{document}
\maketitle

\noindent\emph{Source and licence.} Cavity shape reconstruction with a homogeneous Robin condition via a constrained coupled complex boundary method with ADMM. Version arXiv:2605.12202v2 (CC BY 4.0), distributed by arXiv under the Creative Commons Attribution 4.0 International licence, http://creativecommons.org/licenses/by/4.0/, as recorded in the arXiv licence metadata for this exact version and shown on the abstract page https://arxiv.org/abs/2605.12202v2. Full licence notice: \texttt{licenses/arxiv-2605.12202-CC-BY-4.0.txt}.\par\smallskip

\noindent\textit{Editorial scope.} Counted unit: the article's lemma lem:holder\_continuity (source0.tex lines 557--562). The article's complete original proof of it is delivered with it (source0.tex lines 563--583, one \texttt{proof} environment of 21 lines). No mathematics is written, repaired or re-derived: the statement and the proof are the article's own lines, in the article's own notation, and no retained line is substituted or re-flowed.  Retained uncounted as the source-local context the proof needs: lines 435--439; lines 539--545; lines 551--553.  Nothing was omitted from any retained span, so the package carries no omission record.  No retained line contains a citation, so the wrapper carries no bibliography and no bibliography entry.  No retained line contains a figure environment, an image-inclusion command or drawing code, so no image is delivered and the package declares no asset dependency.  Editorial formatting only, in addition: an AMS \texttt{amsart} preamble that restates the article's own declarations and macros used by the retained lines, namely the environments \texttt{lemma} on separate counters, together with the standard packages those lines need.  The retained environments print in this document as Lemma 1 in order of appearance, whereas the article prints the same items with the numbering of its own full document.  The complete native source of the article as distributed in the arXiv e-print for this exact version is bundled unchanged as \texttt{sources/200407/source0.tex}.
\begin{equation}\label{eq:continuity_of_the_maps}
[t \mapsto {\Jact}] \in {C}^1(\intervalI,{C}(\overline{\varOmega})), \qquad
[t \mapsto A_t] \in {C}^1(\intervalI,{C}(\overline{\varOmega})^{d\times d}),\qquad
[t \mapsto B_t] \in {C}^1(\intervalI,{C}(\partial{\varOmega}))
\end{equation}

\begin{lemma}\label{lem:sesquiliner_form}
	The sesquilinear form $\aat$ defined on $\HH^{1}(\varOmega) \times \HH^{1}(\varOmega)$ by
	\[
		\aat({u},{v}) \coloneqq \intO{A_t \nabla {u} \cdot \nabla {\cv}} + i \intS{{u} {\cv}} + \alpha \intG{B_t {u} {\cv}}, \quad \forall {u}, {v} \in \HH^{1}(\varOmega),
	\]
	is bounded and coercive on $\HH^{1}(\varOmega) \times \HH^{1}(\varOmega)$ for $t \in \intervalI$.
\end{lemma} 

	\begin{equation}\label{eq:transformed_state}
		\aat({u^{t}},{v}) = l({v}), \quad \forall {v} \in \HH^{1}(\varOmega).	
	\end{equation}

\begin{lemma}\label{lem:holder_continuity}
	The solution ${u^{t}}$ of \eqref{eq:transformed_state} is (uniformly) bounded in the neighborhood of $0$ and 
	\[
		\vertiii{u^{t} - u}_{\HH^{1}(\varOmega)} = O(t).
	\] 
\end{lemma}

\begin{proof}
	It is clear that, with $v = u^{t}$ in \eqref{eq:transformed_state} and by Lemma \ref{lem:sesquiliner_form}, $\vertiii{u^{t}}_{\HH^{1}(\varOmega)}$ is bounded for all $t\in \intervalI$.
	Meanwhile, the variational equation $\aat(u^{t},v) - \aaa(u,v) = 0$, for all ${v} \in \HH^{1}(\varOmega)$, is equivalent to
	%
%%%	\begin{equation*}
%%%	\begin{aligned}
%%%	&\intO{\nabla {(u^{t} - u)} \cdot \nabla {\cv}}  + i \intS{(u^{t} - u) {\cv} } + \alpha \intG{(u^{t} - u) {\cv} }\\
%%%		&\qquad = - \intO{(A_{t} - \Id) \nabla {u^{t}} \cdot \nabla {\cv}} - i \intS{ (B_{t} - 1) {u^{t}} {\cv}} - \alpha \intG{ (B_{t} - 1) {u^{t}} {\cv}},
%%%	\end{aligned}
%%%	\end{equation*}
	\begin{equation}\label{eq:difference_equation}
%%%		\aaa(u^{t} - u, v) = - \intO{(A_{t} - \Id) \nabla {u^{t}} \cdot \nabla {\cv}} - i \intS{ (B_{t} - 1) {u^{t}} {\cv}} - \alpha \intG{ (B_{t} - 1) {u^{t}} {\cv}},
	\aaa(u^{t} - u, v) = - \intO{(A_{t} - \Id) \nabla {u^{t}} \cdot \nabla {\cv}} - \alpha \intG{ (B_{t} - 1) {u^{t}} {\cv}},\quad \forall {v} \in \HH^{1}(\varOmega).
	\end{equation}
	% 
	By taking $v = u^{t} - u$ above, and then applying the continuity of the maps given in \eqref{eq:continuity_of_the_maps}, it can be shown without difficulty that the following limit holds
	\[
		\lim_{t \searrow 0} \frac{1}{t} \vertiii{u^{t} - u}_{\HH^{1}(\varOmega)}^{2} = 0,	
	\]
	which verifies the assertion.
\end{proof}

\end{document}
